\section{Conclusions}

The proposed fifth ARC-ADHD paper treats Muse S Athena EEG as an optional measurement instrument rather than a therapeutic mechanism. The primary scientific target is reproducible within-child neurophysiological dynamics around controlled challenge and monitoring-anchored recovery.

The study is falsifiable at several levels: acquisition may be insufficiently synchronized, usable EEG coverage may be too low, spectral effects may be dominated by movement, neural recovery may fail to reproduce, or EEG may add no information beyond behavior. Any of these outcomes is sufficient to reject the need for EEG in the core ARC-ADHD measurement architecture.

Before implementation, the project must obtain and version-lock Athena-specific research acquisition documentation, authorized research access, channel/montage metadata, timestamp semantics, signal-quality rules, synchronization tolerance, and the final preprocessing/analysis specification.
